\section{Dynamic Token Embeddings (EELBERT)}
\label{sec:annex-dynamic-embeddings}

\subsection{Executive Summary}

\texttt{embedding.dynamic} replaces the trainable $V\times H$ input embedding table with an on-the-fly \emph{$n$-gram pooling hash} of the token string, following EELBERT \citep{cohn_eelbert_2023}. The lookup table disappears from the parameter count, from checkpoints and from native deployment artifacts. The effect is largest in the small models where the lookup layer dominates the budget---in BERT-tiny the input embedding accounts for roughly $90\%$ of the parameters---and it is complementary to quantization, pruning and distillation.

Dynamic embeddings are \emph{frozen}: they are never trained and never receive gradients. They only change how a token id becomes a vector; everything downstream (mixers, normalization, residuals, routing, optimizers) is unchanged. As in the paper, the output head remains an independent, untied matrix: during masked-LM pre-training it still stores a $V\times H$ matrix, so the size advantage materialises in the backbone-only or post-fine-tuning artifact (the training checkpoint is not automatically $15\times$ smaller). The implementation lives in \path{src/model/embeddings/dynamic_embedding.py}, the configuration contract is the \texttt{model.embedding.dynamic.*} block, and the tokenizer vocabulary is injected by the training, deployment and inference pipelines through \path{src/utils/tokenizer_vocab.py}.

\subsection{Algorithm}

Let $w$ be a token string of $l$ bytes with byte values $b_1,\dots,b_l$, and let $d=H$ be the embedding width. The algorithm has five steps, with $N$ the maximum $n$-gram size and $B$ the hash bucket.

\begin{enumerate}
\item \textbf{Seeds.} Draw $d$ seeds from a single stored random state $s$,
\[
\mathbf{h} = (h_1,\dots,h_d) \in \{0,\dots,B-1\}^d,
\qquad \mathbf{h} \sim \mathrm{Uniform}(\text{seed}=s).
\]
Only the integer $s$ is persisted; the seed vector is a deterministic function of it.

\item \textbf{$i$-gram signatures.} For $i=1,\dots,N$ compute the polynomial rolling hash of every $i$-gram,
\[
\sigma_i(j) = \Big(\sum_{t=0}^{i-1} b_{j+t}\,a^{\,i-1-t}\Big) \bmod B,
\qquad j=1,\dots,k_i,
\qquad k_i=\max(l-i+1,0),
\]
with $a=257$ (prime, larger than any UTF-8 byte). The hashes are computed incrementally: the $(i{+}1)$-gram hash reuses the $i$-gram hash and appends one byte, so the whole token costs $\mathcal{O}(lN)$ integer operations.

\item \textbf{Partitioned projection.} Split $\mathbf{h}$ into $N$ contiguous parts $\mathbf{h}_i$ of length $d_i$ with
\[
\sum_{i=1}^{N} d_i = d, \qquad d_i \propto i,
\]
so longer $n$-grams receive more dimensions. The projection matrix of each gram size is the outer product
\[
\mathbf{P}_i = \sigma_i \otimes \mathbf{h}_i \in \mathbb{Z}^{k_i\times d_i}.
\]

\item \textbf{Folding.} Map each entry into $[-1,1]$ through
\[
\phi(x) = \frac{(x \bmod B) - \mathbf{1}[x \bmod B > B/2]\;B}{B/2} \in (-1, 1],
\]
which keeps the embedding bounded without storing or normalising a table.

\item \textbf{Pooling and concatenation.} Average each projection across its $k_i$ rows and concatenate the $N$ pieces:
\[
\mathbf{e}_i = \frac{1}{k_i}\sum_{j=1}^{k_i} \phi\big(\mathbf{P}_i[j,:]\big),
\qquad
\mathbf{e} = [\mathbf{e}_1;\dots;\mathbf{e}_N] \in \mathbb{R}^{d}.
\]
When a token is shorter than $i$, $k_i=0$ and that partition contributes a zero vector.
\end{enumerate}

Because the hash is polynomial over characters, morphologically related tokens (e.g.\ \texttt{runner} and \texttt{running}) share many $1$/$2$/$3$-grams and land close in embedding space. The paper shows this matters most for small models: replacing the $n$-gram hash with a purely random (position-preserving) hash costs $5.0$ GLUE points for BERT-mini but only $0.3$ for BERT-base.

\subsection{Configuration Contract}

\begin{table}[H]
\centering
\small
\begin{tabularx}{\linewidth}{@{}l c X@{}}
\toprule
\textbf{Key (\texttt{model.embedding.dynamic.*})} & \textbf{Default} & \textbf{Effect} \\
\midrule
\texttt{enabled} & \texttt{false} & Replace the input lookup table with the frozen hash embedding. Mutually exclusive with \texttt{embedding.factorized.enabled}. \\
\texttt{ngram\_max} & \texttt{3} & Maximum $n$-gram size $N$; must satisfy $1 \le N \le H$. \\
\texttt{bucket} & \texttt{1000000007} & Hash bucket $B$ ($10^9+7$ in the paper). Changing it changes every embedding. \\
\texttt{seed} & \texttt{42} & Random state $s$ for the seeds; identical seeds reproduce identical embeddings. \\
\texttt{precompute} & \texttt{true} & Memoise the $V\times H$ table in a \emph{non-persistent} buffer (pure caching; never saved or exported). \\
\bottomrule
\end{tabularx}
\caption{Dynamic-embedding configuration. The block composes with \texttt{embedding.conv.*}, which applies a Conv1d over the $H$-wide dynamic stream.}
\label{tab:dynamic-embedding-config}
\end{table}

The schema enforces the mutual exclusion with factorized embeddings twice: a conditional \texttt{if/then} rule in \texttt{src/schema/\_conditional\_rules.yaml} and a \texttt{\_\_post\_init\_\_} check in \texttt{FrankensteinModelConfig}. The block is only valid for the text classes (\texttt{frankenstein\_text\_encoder} and \texttt{frankenstein\_text\_decoder}); the ViT uses a convolutional patch embedding and the audio frontend is not a vocabulary lookup, so \texttt{config\_loader} rejects the combination for those classes.

\subsection{Vocabulary Injection and Lifecycle}

Hash embeddings consume token \emph{strings}, while models receive integer ids. The configuration therefore carries a runtime-only, non-schema field \texttt{token\_vocab} (id-ordered token strings). \texttt{inject\_token\_vocab(config, tokenizer)} extracts it from a HuggingFace tokenizer (\texttt{get\_vocab()}) or from the legacy SentencePiece wrapper (\texttt{vocab}, token $\to$ id), attaches it before \texttt{build\_model}, and aligns \texttt{vocab\_size} with the tokenizer ids. The training engine (\path{src/engine.py} and \path{src/training/main.py}), deployment checkpoint loader and inference engine all call the helper, so the vocabulary need not appear in the YAML. Constructing \texttt{DynamicEmbedding} without a vocabulary raises an actionable error.

The checkpoint stores no embedding weights: \texttt{state\_dict()} contains nothing for the input layer, and the vocabulary travels with the tokenizer saved next to the checkpoint. With \texttt{precompute: true} the computed table is a non-persistent buffer---it follows \texttt{.to(device)} but is excluded from checkpoints and exports, so it is exactly equivalent to the on-the-fly computation. With \texttt{precompute: false} the hash is recomputed per batch (lower resident memory, slower steps); the two paths agree up to float32 reduction order.

The native deployment artifacts and the GGUF exporter therefore keep the size benefit. The HuggingFace Transformers export cannot run the custom hash, so it materialises the computed table as a frozen \texttt{nn.Embedding}, records the conversion in the compatibility report and emits a warning that the EELBERT size advantage is not preserved in that artifact.

\subsection{Size Accounting and Paper Results}

The removed block is $V\times H$ weights; the added block is only the optional Conv1d ($k H^2 + H$) when \texttt{embedding.conv.enabled} is true. In our preset \path{configs/eelbert_tiny.yaml} (BERT-tiny geometry: 2 layers, $H=128$, 2 heads, FFN reduced to $128$, $V=30{,}522$), the input table removes $3.91$M weights while the untied MLM head keeps $3.94$M, leaving $4.15$M total parameters---the paper's own accounting, where the head is dropped only for deployment.

Table~\ref{tab:eelbert-results} summarises the reference results reported by \citet{cohn_eelbert_2023}.

\begin{table}[H]
\centering
\small
\begin{tabular}{@{}l r r r r r@{}}
\toprule
\textbf{Model} & \textbf{Params} & \textbf{Size} & \textbf{GLUE} & \textbf{$\Delta$GLUE} & \textbf{Latency} \\
\midrule
BERT-base & 109.5M & 438 MB & 0.775 & --- & 162.0 ms \\
EELBERT-base & 86.1M & 344 MB & 0.760 & $-1.5$ & 165.0 ms \\
BERT-mini & 11.17M & 44.8 MB & 0.753 & --- & 7.0 ms \\
EELBERT-mini & 3.36M & 13.4 MB & 0.746 & $-0.7$ & 9.9 ms \\
BERT-tiny & 4.39M & 17.7 MB & 0.671 & --- & 1.7 ms \\
EELBERT-tiny & 0.48M & 2.04 MB & 0.666 & $-0.5$ & 3.9 ms \\
UNO-EELBERT & 0.31M & 1.24 MB & 0.632 & $-3.9$ & --- \\
\bottomrule
\end{tabular}
\caption{Reference EELBERT results (Table~2 of \citealp{cohn_eelbert_2023}); latency on an M1 MacBook Pro with 32\,GB of RAM. UNO-EELBERT is EELBERT-tiny with the intermediate size reduced to $128$.}
\label{tab:eelbert-results}
\end{table}

The complexity of a token embedding is $\mathcal{O}(l\,d)$ against the $\mathcal{O}(1)$ table lookup of BERT: for the 128--768 widths used by BERT the additional latency is negligible at the base scale but becomes visible on tiny models (about $2.3\times$ for EELBERT-tiny in the table). The paper proposes three mitigations that this implementation leaves as engineering options: a sparse fixed-width mask on the outer product ($\mathcal{O}(l\,s)$ with $s\ll d$), parallel per-token hashing, and a bounded cache of the most frequent $n$-grams.

\subsection{Implementation Map}

\begin{itemize}[leftmargin=1.4em]
\item \path{src/model/embeddings/dynamic_embedding.py} --- \texttt{DynamicEmbedding}: seeds, chunked rolling hash, partitioned outer product, folding, pooling, optional Conv1d, \texttt{embedding\_table()}.
\item \path{src/model/frankenstein_encoder.py} --- \texttt{emb} dispatch (\texttt{dynamic} $\to$ \texttt{factorized} $\to$ \texttt{nn.Embedding}); the decoder reuses it through its encoder backbone. Mutual exclusion and range checks live in \path{src/model/config.py}.
\item \path{src/utils/tokenizer_vocab.py} --- \texttt{token\_strings()} and \texttt{inject\_token\_vocab()}, shared by the engine, the CLI training path, deployment and inference.
\item \path{src/deploy/transformers\_export.py} --- frozen-table materialisation and the export warning.
\item \path{tests/test\_dynamic\_embedding.py} --- determinism per seed, $[-1,1]$ bounds, morphology similarity, precompute vs.\ on-the-fly parity, zero trainable parameters, vocabulary and configuration errors, encoder/decoder forwards, schema rule and export materialisation.
\item \path{configs/eelbert_tiny.yaml} and \path{configs/examples/embedding_dynamic.yaml} --- presets; both are smoke-tested, and the E2E harness covers the combination as \texttt{trans\_dynamic\_emb}.
\end{itemize}

\subsection{Scope and Limitations}

Dynamic embeddings are a compression mechanism, not a new representation-learning objective: they do not approximate a trained BERT embedding, and they cannot be trained or fine-tuned themselves. Because the mapping is a deterministic function of (string, seed, $B$, $N$), changing any of those invalidates previously trained checkpoints. The design also couples the model to the tokenizer that produced the training ids: the vocabulary (and therefore the token strings) is part of the artifact contract, and a tokenizer swap silently changes every input embedding. Finally, the HuggingFace export loses the size advantage by construction, so on-device deployments should prefer the native checkpoint, the deployment package or GGUF.
